% ID: SEL-D0040-OB-E02
% ESTADO: validada. S041, cotejada con examen y criterios.
\textbf{a)} $f'(x)=(3x-1)(x-1)$ es positiva fuera de
$[1/3,1]$ y negativa dentro. Por tanto hay máximo relativo
en $\boxed{(1/3,4/27)}$ y mínimo relativo en $\boxed{(1,0)}$.
La segunda derivada $6x-4$ cambia de signo en $2/3$.
Como $f(2/3)=2/27$, la inflexión es $\boxed{(2/3,2/27)}$.

\textbf{b)} $f(x)=x(x-1)^2$. Corta el eje horizontal
en $(0,0)$ y $(1,0)$, y el vertical en $(0,0)$.

\textbf{c)} En $0$ la pendiente es 1 y en $1$ es 0.
Como ambos puntos están en el eje horizontal,
las tangentes son $\boxed{y=x}$ y $\boxed{y=0}$, respectivamente.
